\section{System Model and Design Overview}
\label{sec:system-model}

We consider two unbounded input streams, $L$ and $R$, where each tuple is an embedding vector with metadata $(\mathit{uid}, \mathit{ts}, \mathbf{v})$. The task is a sliding-window similarity join: for each arriving tuple, the system finds tuples from the opposite stream that are both temporally valid and semantically similar. Let $W$ be the window size and $\theta$ be the similarity threshold. A pair $(l, r)$ is emitted iff $|l.\mathit{ts} - r.\mathit{ts}| \leq W$ and $\mathrm{sim}(l, r) \geq \theta$. In SageFlow, similarity is computed with configurable policies (e.g., fixed or adaptive $\alpha$ in exponential distance-to-similarity mapping), while execution uses eager processing so that each arrival immediately updates state and probes the opposite side.

The main systems challenge is not candidate search alone, but maintaining correctness and throughput under concurrent window evolution. In multicore streaming execution, tuples are inserted, expired, and queried continuously; therefore, partitioning, state layout, and index policy must be co-designed. A key implementation lesson in our codebase is that invalid combinations can silently hurt recall even when per-operator throughput appears high. For example, RoundRobin routing with non-shared window state fragments logically related tuples across workers and may miss cross-partition matches. This motivates explicit compatibility rules between partition strategy and state/index strategy as part of the runtime contract.

SageFlow addresses this with a strategy-driven join architecture. A join pipeline is instantiated by four coupled dimensions: (i) join algorithm (BruteForce, IVF, HNSW, HDR-Tree, LSH, ClusteredJoin, S3J-style, and VSJoin-oriented path), (ii) partition strategy (RoundRobin, key/vector hash, LSH, centroid), (iii) window-state type (shared or partitioned variants), and (iv) index strategy (shared or partitioned). This decomposition enables a single execution framework to host multiple baselines and research ideas while reusing common components such as RuntimeContext, WindowState abstraction, ConcurrencyManager-mediated index access, and unified output verification.

Our VSJoin-oriented design follows two principles. First, it separates write-efficient local ingestion from read-efficient global recall by combining partition-local indexes with a global index layer. Local indexes reduce write-path contention and preserve partition locality, while global indexes increase cross-partition candidate coverage. Second, it separates candidate generation from final verification: index probes produce candidate sets, and temporal/similarity predicates validate output pairs. This split keeps the verification semantics uniform across methods and makes the quality-latency trade-off explicit at the candidate stage.

To preserve join quality near partition boundaries, routing is boundary-aware and can multicast records to multiple candidate partitions. In addition, the runtime supports periodic global-index refresh to align read structures with evolving window state. These mechanisms are intentionally lightweight: they avoid heavyweight global synchronization in the hot path while providing practical guardrails against recall degradation in high-parallel settings. For reproducibility, we keep algorithm-specific constraints explicit (e.g., compatible partition/state pairs, partition-count alignment for centroid-partitioned modes) and evaluate all methods within the same end-to-end harness.

This model and design provide the foundation for the detailed algorithms and implementation described in subsequent sections. In particular, the rest of the paper focuses on how candidate generation, routing, and state/index maintenance interact under streaming windows, and why these interactions dominate real multicore behavior beyond standalone ANN query benchmarks.
